\documentclass[11pt]{article}
\usepackage{amsmath,amssymb,amsthm,mathtools}

\newtheorem{theorem}{Theorem}
\newtheorem{lemma}[theorem]{Lemma}
\newtheorem{corollary}[theorem]{Corollary}
\newcommand{\rank}{\operatorname{rank}}
\newcommand{\im}{\operatorname{im}}

\title{An Obstruction to a Universal Fermionic Statistics-Carrier Tensor Factorization}
\author{Working theorem draft}
\date{2026-09-01}

\begin{document}
\maketitle

\paragraph{Evidence status.}
The dimension obstruction is an elementary proof. An implementation-independent
exact-rational verifier checks all exterior flattening ranks for
$3\le N\le8$. Its certificate hash is
\begin{center}
\small\texttt{06025f168a49c0ab857c2163103ffabcb56fb04cd1fed4df9120d25ef6bc60df}.
\end{center}
The representation-theory statements use standard sign-representation and
Littlewood--Richardson facts. External proof review is still required.

\section{Cut object and proposed factorization}

Let $0\ne C\in\Lambda^N V$. At a particle cut $k\mid N-k$, define
\[
 c_{C,k}:\Lambda^kV^*\longrightarrow\Lambda^{N-k}V,
 \qquad c_{C,k}(\alpha)=\iota_\alpha C.
\]
Its rank is the particle-cut Schmidt rank, up to a fixed normalization of the
exterior coproduct. Write $\mathcal B_k(C)=\im c_{C,k}$.

The literal statistics-carrier proposal asks for finite-dimensional spaces
with
\begin{equation}
 \mathcal B_k(C)\cong
 \mathcal S^{\mathrm{fermion}}_{N,k}\otimes\mathcal M_k(C),
 \label{eq:factor}
\end{equation}
where the first factor is fixed by $(N,k)$ and every nonzero Slater determinant
has $\dim\mathcal M_k=1$.

\section{Dimension obstruction}

\begin{theorem}
For every $N\ge3$, no factorization of the form \eqref{eq:factor} is exact for
all nonzero $N$-forms.
\end{theorem}

\begin{proof}
At the $1\mid N-1$ cut, every nonzero Slater determinant has rank $N$.
The multiplicity-one condition therefore forces
$\dim\mathcal S^{\mathrm{fermion}}_{N,1}=N$, so every rank admitted by
\eqref{eq:factor} must be divisible by $N$.

Let $\dim V=N+2$ and, with the empty product interpreted as one, set
\[
 W=e_6\wedge\cdots\wedge e_{N+2},\qquad
 C_N=e_1\wedge(e_2\wedge e_3+e_4\wedge e_5)\wedge W.
\]
The $N+2$ contractions $\iota_{e_i^*}C_N$ are linearly independent.
For $i=2,3,4,5$ each contains a distinct monomial with the full factor $W$.
The $i=1$ contraction is the only one without $e_1$. For $i\ge6$, the two
monomials have a distinct missing factor of $W$. These monomial supports cannot
cancel across the displayed groups. Hence
\[
 \rank c_{C_N,1}=N+2.
\]
But $N\nmid N+2$ for $N\ge3$, contradicting \eqref{eq:factor}.
\end{proof}

The counterexample is a sum of two Slaters. If the two component contraction
maps were supplied independent covectors, their images would give $2N$
dimensions. In the physical sum the same covector acts on both maps, locking
$N-2$ shared-orbital channels and leaving rank $N+2$. Thus a free
two-dimensional multiplicity misses state-dependent linear relations even at
Slater-sum length two.

\section{What symmetry adaptation actually gives}

Symmetry-adapted tensor networks split invariant tensors into fixed
group-theoretic structural tensors and free degeneracy tensors
\cite{SinghPfeiferVidal2010,Weichselbaum2012}. The relevant symmetries here do
not produce \eqref{eq:factor}.

\begin{lemma}
Particle-permutation symmetry supplies a one-dimensional sign carrier. Under
$GL(V)$, the exterior coproduct supplies a multiplicity-one equivariant
coupling, while retaining the full orbital exterior representation.
\end{lemma}

\begin{proof}
The antisymmetric sector is the one-dimensional sign representation of $S_N$.
On restriction to $S_k\times S_{N-k}$ it is
$\operatorname{sgn}_k\otimes\operatorname{sgn}_{N-k}$, again once. Therefore
$\binom Nk$ is not an $S_N$ structural-irrep or degeneracy dimension.

For $GL(V)$, $\Lambda^N V$ is irreducible. The Littlewood--Richardson
coefficient of $\Lambda^N V$ in
$\Lambda^kV\otimes\Lambda^{N-k}V$ is one, so the exterior coproduct is the
unique equivariant coupling up to scale. Treating $C$ as a covariant external
irrep vector retains $\binom DN$ orbital components; it does not isolate a
smaller correlation multiplicity.
\end{proof}

Hamiltonian-specific orbital symmetry groups can yield useful multiplet and
charge degeneracies. Those are established symmetry-adapted TN structures,
depend on the physical symmetry and sector, and do not universally give Slater
multiplicity one.

\section{State-adaptive alternatives}

For a decomposable form, the occupied $N$-plane is intrinsic and its exterior
powers explain the binomial cut rank. A general correlated form has no
canonical occupied $N$-plane. Choosing a Slater/secant expansion instead gives
an established determinant expansion. Slater rank is standard already for two
fermions \cite{SchliemannEtAl2001}; in higher degree uniqueness becomes a
Grassmannian-secant identifiability problem with nonuniform geometry
\cite{BallicoEtAl2013,GalganoStaffolani2024}.

The ordinary singular values of $c_{C,k}$ are canonical under orbital
unitaries and give the usual truncation spectrum. They reproduce the flat
$\binom Nk$ Slater spectrum, so
they do not remove the exchange rank. A chosen nonorthogonal Slater/AGP sum can
be truncated only with its overlap Gram matrix and an explicit state or energy
error bound; it is not a canonical tensor factor of $C$.

\section{The compact-input contraction cost remains}

\begin{corollary}
An intrinsic multiplicity-one label for Slater states does not imply a
polynomial compiler or norm contraction from compact matrix-wedge input.
\end{corollary}

\begin{proof}
The sparse-path construction of the companion contraction theorem produces
\[
 \Psi_A=\frac{\operatorname{perm}(A)}{M!}
 e_1\wedge\cdots\wedge e_{2M}.
\]
Whenever nonzero, this is projectively a single Slater and therefore must have
correlation multiplicity one under the proposed sanity condition. Yet obtaining
its scale and squared norm from the compact path data computes the permanent.
Thus the hard scalar lies outside the state-intrinsic label.
\end{proof}

This corollary concerns the required norm and reconstruction algorithm. It does
not claim that every normalized observable on this one-dimensional output
family is hard.

\paragraph{Decision.}
The direct universal statistics-carrier tensor product fails exactness,
canonicality, and compact-input contraction. Symmetry-adapted charge networks
and state-adaptive finite Slater/AGP expansions remain valid established
methods, not the missing universal FEMPS multiplicity spectrum.

\begin{thebibliography}{9}

\bibitem{SinghPfeiferVidal2010}
S.~Singh, R.~N.~C.~Pfeifer, and G.~Vidal,
``Tensor network decompositions in the presence of a global symmetry,''
\emph{Physical Review A} \textbf{82}, 050301 (2010).

\bibitem{Weichselbaum2012}
A.~Weichselbaum,
``Non-Abelian symmetries in tensor networks: a quantum symmetry space
approach,'' \emph{Annals of Physics} \textbf{327}, 2972--3047 (2012).

\bibitem{SchliemannEtAl2001}
J.~Schliemann, J.~I.~Cirac, M.~Ku\'s, M.~Lewenstein, and D.~Loss,
``Quantum correlations in two-fermion systems,''
\emph{Physical Review A} \textbf{64}, 022303 (2001).

\bibitem{BallicoEtAl2013}
E.~Ballico, A.~Bernardi, M.~V.~Catalisano, and L.~Chiantini,
``Grassmann secants, identifiability, and linear systems of tensors,''
\emph{Linear Algebra and its Applications} \textbf{438}, 121--135 (2013).

\bibitem{GalganoStaffolani2024}
V.~Galgano and R.~Staffolani,
``Identifiability and singular locus of secant varieties to Grassmannians,''
\emph{Collectanea Mathematica} (2024).

\end{thebibliography}

\end{document}
